\section{Jeu de données}

\subsection{Vue d'ensemble}

Nous utilisons un jeu de données IRM de tumeurs cérébrales publiquement disponible,
contenant \NTotal{} images réparties équitablement entre \NClasses{} classes :
\textbf{gliome}, \textbf{méningiome} et \textbf{tumeur hypophysaire}
(Tableau~\ref{tab:dataset}).

\begin{table}[H]
  \centering
  \caption{Synthèse du jeu de données.}
  \label{tab:dataset}
  \begin{tabular}{@{}lc@{}}
    \toprule
    \textbf{Propriété} & \textbf{Valeur} \\
    \midrule
    Nombre de classes & \NClasses{} (gliome, méningiome, hypophysaire) \\
    Nombre total d'images & \NTotal{} \\
    Échantillons par classe & \NPerClass{} \\
    Partition d'entraînement & \NTrain{} (80\,\%, stratifiée) \\
    Partition de test & \NTest{} (20\,\%, stratifiée) \\
    Échantillons de test par classe & \NTestPerClass{} \\
    Taille d'image (après prétraitement) & \ImgSize{} niveaux de gris \\
    Graine aléatoire & 42 \\
    \bottomrule
  \end{tabular}
\end{table}

\subsection{Équilibre des classes}

Le jeu de données est parfaitement équilibré : chaque classe contribue de manière
égale aux partitions d'entraînement et de test. Cette conception élimine le biais
lié au déséquilibre des classes et garantit que les métriques macro-moyennées
reflètent directement la capacité de discrimination par classe, sans être biaisées
par une classe dominante (Figure~\ref{fig:class_balance}).

\begin{figure}[H]
  \centering
  \includegraphics[width=0.55\linewidth]{class_balance_train.png}
  \caption{Répartition des classes dans la partition d'entraînement. Chaque classe
           compte 400 échantillons d'entraînement.}
  \label{fig:class_balance}
\end{figure}

\subsection{Choix des métriques d'évaluation}

Compte tenu de l'équilibre des classes, nous utilisons le \textbf{F1-macro} comme
métrique principale de classement. Elle accorde un poids égal aux trois types
tumoraux et mesure directement la capacité du pipeline à distinguer chaque classe.
La précision (accuracy) et le $\kappa$ de Cohen sont rapportés comme métriques secondaires.
